\documentclass[11pt]{article}
\usepackage[margin=2.5cm]{geometry}
\usepackage{xcolor}
\usepackage{booktabs}
\usepackage{array}
\setlength{\parskip}{6pt}
\setlength{\parindent}{0pt}
\usepackage{titlesec}
\usepackage{fancyhdr}
\usepackage{fontenc}
\usepackage{inputenc}

% ── Palette ───────────────────────────────────────────────────────────────────
\definecolor{accent}{HTML}{1A73E8}
\definecolor{accent2}{HTML}{18965C}
\definecolor{lightgray}{HTML}{555555}
\definecolor{codefill}{HTML}{F1F3F4}
\definecolor{rulegray}{HTML}{CCCCCC}

% ── Section styling ───────────────────────────────────────────────────────────
\titleformat{\section}{\large\bfseries\color{accent}}{}{0em}{}[\color{rulegray}\rule{\linewidth}{0.6pt}]
\titleformat{\subsection}{\normalsize\bfseries\color{accent2}}{}{0em}{}
\titlespacing{\section}{0pt}{16pt}{6pt}
\titlespacing{\subsection}{0pt}{10pt}{3pt}

% ── Code/formula box ──────────────────────────────────────────────────────────
\newenvironment{codebox}{%
  \vspace{6pt}\par
  \noindent\hspace{1cm}%
}{%
  \par\vspace{6pt}%
}

% ── Header / footer ───────────────────────────────────────────────────────────
\pagestyle{fancy}
\fancyhf{}
\renewcommand{\headrulewidth}{0.4pt}
\fancyhead[L]{\small\color{lightgray}Pipeline Notes}
\fancyhead[R]{\small\color{lightgray}May 7, 2026}
\fancyfoot[C]{\small\color{lightgray}\thepage}

\begin{document}

% ── Title ─────────────────────────────────────────────────────────────────────
{\LARGE\bfseries\color{accent} Pipeline Notes}\par
{\large\color{lightgray} May 7, 2026 \ $\cdot$\ Daniel Chang\ $\cdot$\ Franco Lab}
\vspace{4pt}
\color{rulegray}\rule{\linewidth}{1.2pt}\color{black}
\vspace{8pt}

% ── Big picture ───────────────────────────────────────────────────────────────
\section*{The Big Picture}

The pipeline answers one question: \textbf{how much more concentrated is a protein inside a condensate compared to the surrounding nucleus?} That ratio is the Partition Coefficient (PC). To compute it you need to know \textit{where} the condensates are and \textit{where} the nucleus is --- that's what the pipeline figures out automatically.

% ── Steps ─────────────────────────────────────────────────────────────────────
\section*{Step-by-Step}

\subsection*{1.\ \ Load}
Two 3D image stacks are read in from TIF files:
\begin{itemize}
  \item \textbf{Ch1} --- the nuclei channel (fluorescent signal from the nucleus)
  \item \textbf{Ch2} --- the condensate channel (fluorescent signal from the protein of interest)
\end{itemize}
Each stack is 55 Z-slices $\times$ 185 $\times$ 259 pixels. Think of it as a 3D block of pixels, like a loaf of bread sliced 55 times.

\subsection*{2.\ \ Denoise}
Before trying to find any boundaries, every Z-slice gets run through Cellpose's denoising model. This sharpens edges and suppresses shot noise so the segmentation model can draw tighter, more accurate boundaries.

\textbf{Key point:} the denoised image is thrown away after this step. It is only used to guide \textit{where} the boundaries go --- never to measure intensity values.

\subsection*{3.\ \ Segment}
Cellpose 3 (cyto3 model, \texttt{do\_3D=True}) processes the full 3D stack --- not slice by slice, but all 55 slices simultaneously. This means it tracks the same object consistently across Z rather than redrawing boundaries independently on each slice.

The output is two sets of \textbf{masks}:
\begin{itemize}
  \item \textbf{Condensate masks} --- each condensate gets a unique integer label. Every pixel inside condensate \#1 is labeled 1, condensate \#2 is labeled 2, and so on.
  \item \textbf{Nuclei masks} --- same idea for nuclei.
\end{itemize}

After segmentation, the nuclei masks go through one extra step: Cellpose tends to split each nucleus into 15--25 fragments because of internal texture (the condensates inside the nucleus confuse it). To fix this, the binary mask (any pixel labeled $> 0$ = ``inside a nucleus'') is run through 3D connected-component analysis. This collapses all the fragments back into the true connected regions --- 76 Cellpose labels become 5 clean nuclei.

\subsection*{4.\ \ Measure}
Using the masks as selection tools on the original raw image, \texttt{regionprops} extracts per-object measurements for every Z-slice: area, centroid, and mean intensity.

\subsection*{5.\ \ 3D Volume}
Counts how many voxels (3D pixels) each labeled object occupies across all 55 slices. This gives a true 3D volume estimate for each condensate and nucleus rather than just a per-slice area.

\subsection*{6.\ \ Partition Coefficient}

This is where the masks are used to select pixels from the \textbf{raw image} (Ch2, the condensate channel).

\textbf{Background ($B$)} \\
Find the minimum pixel value across the entire 3D stack. This is the camera/autofluorescence floor --- a baseline offset that has nothing to do with the protein. It gets subtracted from every pixel before any density is computed.

\textbf{Condensate density}
\begin{itemize}
  \item Take all pixels inside \textit{both} a condensate mask \textit{and} a nucleus mask
  \item Subtract $B$ from each, clip to 0
  \item Sort by intensity, keep only the top 75\% brightest
  \item Take the mean
\end{itemize}
The top-75\% trim is important: Cellpose draws slightly inflated boundaries that include a dim halo around the true bright core. Trimming the dimmest 25\% of mask pixels recovers approximately what a researcher would have drawn by hand.

\textbf{Dilute density}
\begin{itemize}
  \item Take all pixels inside a nucleus mask but \textit{outside} any condensate mask --- this is the quiet nuclear background
  \item Find all valid $10{\times}10{\times}10$ voxel patches that fit entirely within this region (${\sim}88{,}000$ patches)
  \item Sort by mean intensity, take the 50 lowest-intensity patches
  \item Average them
\end{itemize}
The lowest-50-patch approach mimics how a researcher manually picks a quiet, representative background spot. Using a single random patch was unstable (PC would swing from 4.3 to 6.0 depending on where it landed).

\textbf{PC formula}
\begin{codebox}
\texttt{PC = condensate\_density / dilute\_density}
\end{codebox}
Interpretation: how many times brighter are condensates than the surrounding nuclear background, after removing the camera offset.

% ── Why each step matters ─────────────────────────────────────────────────────
\section*{Why Each Step Matters}

\vspace{4pt}
\begin{tabular}{>{\bfseries}p{4.5cm} p{10cm}}
\toprule
Step & What would break without it \\
\midrule
Denoising & Blurry boundaries $\to$ masks include too much dim halo $\to$ PC underestimated \\[4pt]
do\_3D=True & Slice-by-slice segmentation gives inconsistent boundaries across Z $\to$ noisy intensity measurements \\[4pt]
CC nuclei fix & Fragmented nucleus labels leave gaps $\to$ condensates in gaps excluded from PC \\[4pt]
Background subtraction & PC systematically lower than reference (reference always subtracts $B$) \\[4pt]
Top-75\% trim & Cellpose halos drag cond\_density down $\to$ PC underestimated vs.\ manual reference \\[4pt]
Lowest-50-patch dilute & Single random patch is seed-dependent $\to$ PC swings $\pm$30\% across runs \\
\bottomrule
\end{tabular}

% ── Biological meaning ────────────────────────────────────────────────────────
\section*{What PC Means Biologically}

A PC of 1.0 means the protein is equally distributed inside condensates and in the surrounding nucleus --- no enrichment. A PC of 6.3 (our result, matching the manual reference) means the protein is ${\sim}6\times$ more concentrated inside the condensate than in the nuclear background. Higher PC = tighter, more selective condensate.

% ── Batch metrics ─────────────────────────────────────────────────────────────
\section*{Batch Validation Metrics Explained}

The batch run across 30 JABr cells reports three numbers: $r$, RMSE, and bias.

\subsection*{Pearson $r$ (correlation)}
Measures how well the pipeline's PC values \textit{track} the reference values --- do they go up and down together? $r=1.0$ is perfect, $r=0$ is no relationship. $r=0.735$ means the pipeline ranks cells in roughly the right order (high PC cells are still high, low PC cells are still low), but it is not a tight 1:1 relationship.

\subsection*{RMSE (Root Mean Square Error)}
The average size of the error per cell, in the same units as PC. RMSE$=2.815$ means on average the pipeline is off by ${\sim}2.8$ PC units per cell. Since the reference PCs range from ${\sim}4$ to ${\sim}20$, that is a moderate error --- good for most cells, worse for the outliers.

\subsection*{Bias (+2.8\%)}
The systematic over- or under-estimation averaged across all cells. $+2.8\%$ means the pipeline is very slightly over-estimating PC on average --- essentially zero systematic error. A large positive bias would mean the pipeline always reads high; a large negative bias would mean it always reads low. The top-75\% trim was specifically tuned to get this close to zero (it was $-16\%$ before the fix).

\vspace{8pt}
\begin{tabular}{>{\bfseries}p{4cm} p{2.5cm} p{7.5cm}}
\toprule
Metric & Value & What it tells you \\
\midrule
Pearson $r$  & 0.735          & Does it correlate? (shape of relationship) \\[4pt]
RMSE         & 2.815 PC units & How far off is it on average? (absolute error) \\[4pt]
Bias         & +2.8\%         & Does it consistently over/under-shoot? (systematic offset) \\
\bottomrule
\end{tabular}

\end{document}
